\documentclass[tikz,border=10pt]{standalone}

\usepackage{tikz}
\usetikzlibrary{arrows.meta, positioning}

\begin{document}

\begin{tikzpicture}[
    block/.style = {draw, rectangle, minimum height=1.2cm, minimum width=2.8cm, align=center},
    sum/.style   = {draw, circle, inner sep=0pt, minimum size=6mm},
    arrow/.style = {->, thick}
]

% ===== Knoten =====
\node (w) {$w(s)$};
\node[sum, right=1.6cm of w] (sum) {};

\node[block, right=1.6cm of sum] (R) {Regler $G_R(s)$};
\node[block, right=2.2cm of R] (A) {Stellglied $G_S(s)$};
\node[block, right=2.2cm of A] (G) {Regelstrecke $G_{RS}(s)$};

\node[right=1.8cm of G] (y) {$y(s)$};

\node[block, below=1.6cm of G] (H) {Messglied $G_M(s)$};

\node[above=1.2cm of G] (d) {$z(s)$};

% ===== Vorwärtsweg =====
\draw[arrow] (w) -- node[pos=0.9, above] {$+$} (sum);
\draw[arrow] (sum) -- node[above] {$e(s)$} (R);
\draw[arrow] (R) -- node[above] {$u(s)$} (A);
\draw[arrow] (A) -- node[above] {$u_s(s)$} (G);
\draw[arrow] (G) -- node[above] {} (y);

% ===== Störgröße =====
\draw[arrow] (d) -- (G);

% ===== Rückführung =====
\draw[arrow] (y) |- (H);
\draw[arrow] (H) -| node[pos=0.95, left] {$-$} (sum);

% ===== Beschriftungen =====
%\node at ([yshift=5pt]sum.north) {};
\node[below=0.2cm of H] {$y_M(s)$};
\node[below=0.2cm of sum] {};

\end{tikzpicture}

\end{document}
